\section{Выводы}

В ходе выполнения лабораторной работы были изучены и использованы конкретные системные вызовы операционных систем Windows
(Win32 API) и Linux (POSIX) для работы с процессами и анонимными каналами: \texttt{CreatePipe},
\texttt{CreateProcessW}, \texttt{WaitForSingleObject}, \texttt{pipe}, \texttt{fork}, \texttt{dup2},
\texttt{execv}, \texttt{waitpid}.

Разработана программа из двух взаимодействующих через пайпы процессов в соответствии с ТЗ. Программа выполнена на языке
C по схеме \enquote{единый API + две платформозависимые реализации}, благодаря чему логика
приложений полностью кроссплатформенна, а выбор реализации выполняется на этапе сборки.

Реализовано детализированное логирование всех процессов и системных вызовов.
Исправлены ошибки организации проекта, допущенные при старте работы.
Отлажены краевые случаи падения процессов с появлением фантомных исходов или полным их отсутствием.
Достигнута равная работоспособность на Windows и Linux из единой исходной базы.
Подтверждена корректная обработка ошибок и завершение обоих процессов с ненулевым кодом при сбоях.
Сформирован полноценный отчет о ходе и результатах проделанной работы.
\pagebreak